\documentclass[11pt]{article}
\usepackage[margin=1.03in]{geometry}
\usepackage{amsmath,amssymb,amsthm,microtype}
\usepackage[colorlinks=true,linkcolor=blue,citecolor=blue,urlcolor=blue]{hyperref}
\newtheorem{theorem}{Theorem}
\newtheorem{corollary}[theorem]{Corollary}
\newcommand{\T}{\mathbb T}
\newcommand{\R}{\mathbb R}
\newcommand{\E}{\mathbb E}
\newcommand{\Prob}{\mathbb P}
\newcommand{\one}{\mathbf 1}
\title{The sharp universal scale for lacunary averages\\
Erd\H{o}s problem 995}
\author{Taisho \textnormal{(Round 48 research note)}}
\date{24 November 2026}
\begin{document}
\maketitle
\begin{abstract}
For $1<p<\infty$, averages of any fixed $L^p$ function along any sequence of measure-preserving maps are $o((\log N)^{1/p})$ almost surely. This is an optimal universal scale even for dyadic lacunary dilates: we construct one sequence of powers of two such that, for every $p>1$ and every positive gauge $g(N)=o((\log N)^{1/p})$, a mean-zero $f\in L^p$ has $\limsup_N g(N)^{-1}N^{-1}\sum_{j\le N} f(n_j x)=+\infty$ almost everywhere. For nondecreasing gauges we give an exact criterion. In particular the general growth-estimation clause of Erd\H{o}s problem 995 has this sharp universal interpretation, while its specific proposed smaller bound was already refuted by Ho.
\end{abstract}

\section{The question and the endpoint}
For positive integers $n_j$ and a measurable real function $f$ on the circle $\T=\R/\mathbb Z$, write
\[
 S_N(f;x)=\sum_{j=1}^N f(n_jx),\qquad A_N(f;x)=S_N(f;x)/N.
\]
Erd\H{o}s discusses the gap between his lower and upper growth bounds for lacunary $n_j$ in \cite[Section II, p.~58, (15)--(16)]{erdos}. The catalogued \href{https://www.erdosproblems.com/995}{Erd\H{o}s problem 995} asks to estimate the almost-everywhere growth of $S_N$ for square-integrable $f$ and includes the example question whether $S_N=o(N\sqrt{\log\log N})$ universally. The latter assertion has already been disproved by Ho \cite[Theorem 1.2]{ho}. For $2\le p<\infty$ he gives examples beating $N(\log N)^{1/p-\varepsilon}$ for all fixed $\varepsilon>0$, against an upper bound with exponent $1/p+\varepsilon$ in his Remark 6.1. His Section 8 explicitly leaves the order without $\varepsilon$ losses, including secondary logarithmic factors, to be determined. We give a sharp universal-scale answer; this does not assert the same rate for every individual function or every prescribed lacunary sequence.

All unqualified logarithms below are natural. Put $E_p(N)=(\log(N+1))^{1/p}$. A universal estimate means each fixed function and each fixed sequence satisfy it outside a null set, not that one null set works for all choices.

\begin{theorem}[Universal upper bound]\label{thm:upper}
Let $(X,\mu)$ be a probability space, let $T_j:X\to X$ be any measure-preserving maps, and let $f\in L^p(X)$ for $1<p<\infty$. Then
\[
 \frac1N\sum_{j\le N}|f(T_jx)|=o(E_p(N)) \qquad\text{for $\mu$-almost all \(x\)}.
\]
In fact there exist constants $c_p,C_p>0$ such that, for all $M\ge2$ and $u>0$,
\[
 \mu\left\{x:\sup_{N\ge M}\frac{\sum_{j\le N}|f(T_jx)|}{N E_p(N)}>u\right\}
 \le C_p u^{-p}\int_{|f|>c_puE_p(M)}|f|^p\,d\mu .
 \tag{1}\label{eq:max}
\]
\end{theorem}

\begin{theorem}[Universal lower sequence]\label{thm:lower}
There is one sequence of powers of two $n_j=2^{m_j}$, with $n_{j+1}/n_j\to\infty$, independent of $p$ and of $g$, with the following property. For every $1<p<\infty$ and positive $g:\mathbb N\to\R$ with $g=o(E_p)$, some real mean-zero $f\in L^p(\T)$ satisfies
\[
 \limsup_{N\to\infty}\frac{A_N(f;x)}{g(N)}=+\infty
 \qquad \text{for almost all \(x\)}.
 \tag{2}\label{eq:lower}
\]
The same holds for nondecreasing $g$ if merely $\liminf_{N\to\infty} g(N)/E_p(N)=0$.
\end{theorem}

\begin{corollary}
For $1<p<\infty$ and a positive nondecreasing gauge $g$, the bound $A_N(f;x)=o(g(N))$ almost everywhere for every $f\in L^p(\T)$ and every lacunary integer sequence holds if and only if $\liminf g(N)/E_p(N)>0$. In particular the optimal universal scale for problem 995 is $N\sqrt{\log N}$ for the sums, in the sense that this scale gives little-oh and no universally smaller deterministic rate does.
\end{corollary}

\section{Proof of the upper bound}
Let $F=|f|$ and $B_N=N^{-1}\sum_{j\le N} F\circ T_j$. For $J\ge1$ and $v>0$, split $F$ at $vJ^{1/p}/2$. The low part contributes at most this cutoff to $B_{2^J}$. Markov's inequality and measure preservation give
\[
 \mu\{B_{2^J}>vJ^{1/p}\}
 \le \frac{2}{vJ^{1/p}}\int F\one_{\{F>vJ^{1/p}/2\}}\,d\mu .
\]
For $J\ge J_0\ge1$, sum this estimate and use Tonelli:
\[
 \begin{aligned}
 \sum_{J\ge J_0}\mu\{B_{2^J}>vJ^{1/p}\}
 &\le \frac{2}{v}\int F
       \sum_{\substack{J\ge J_0\\J<(2F/v)^p}} J^{-1/p}\,d\mu\\
 &\le C_pv^{-p}\int_{\{F>vJ_0^{1/p}/2\}}F^p\,d\mu .
 \end{aligned}
 \tag{3}\label{eq:sum}
\]
The last estimate uses $\sum_{1\le J<H}J^{-1/p}\le C_p H^{1-1/p}$ for $H\ge1$, which is why we assume $p>1$. For $2^{J-1}<N\le 2^J$, we have $B_N\le2B_{2^J}$ and $E_p(N)\ge c_pJ^{1/p}$. Therefore the event in \eqref{eq:max} is contained in the union of the dyadic events in \eqref{eq:sum} with $J\ge \lceil\log_2 M\rceil$ and $v$ a fixed positive multiple of $u$. The displayed bound follows. First Borel--Cantelli in \eqref{eq:sum}, for positive rational $v$, gives $B_{2^J}=o(J^{1/p})$ almost surely, and the same comparison extends this to all $N$. In particular, no cancellation, lacunarity, or independence is used on the upper side.

\section{One lower sequence for every gauge}
We adapt the independent dyadic-trial mechanism of Ho \cite[Sections 2--4 and 6]{ho}; no Fourier-tail control is required here. Nonnegative binomial blocks and an amplitude choice \emph{after} fixing the sequence make it possible to keep a single sequence for all the gauges.

\subsection*{The fixed sequence and its trial events}
In stage $s\ge1$ take $t_s=2^{s-1}$ trials. Start with zero chosen exponents; whenever the total chosen so far is $P$, append a trial of $P+1$ exponents. The new endpoint is $2P+1$. The total number at the end of stage $s$, denoted $N_s$, satisfies
\[
 N_s+1=2^{t_1+\cdots+t_s}=2^{2t_s-1}.
\]
Set $L_s=N_s+1$ and $d_s=\log_2(L_s t_s)$, a positive integer. Every trial length $\ell$ in stage $s$ satisfies $\ell\le L_s/2$. A trial of length $\ell$ in stage $s$ consists of the exponents
\[
  M,\ M+d_s,\ \ldots,\ M+(\ell-1)d_s.
 \tag{4}\label{eq:exponents}
\]
Choose $M$ deterministically larger than the last exponent and larger than every binary digit position used by the trial events defined below in previous trials. For example initialize $H=0$, take $M=H+1$ and, after the trial, update $H=M+(L_s+1)d_s$. This yields one strictly increasing sequence $m_j$ of exponents. The internal spacing $d_s$ goes to infinity, and a gap after a trial is at least $d_s$ by the update rule. Thus $2^{m_{j+1}-m_j}\to\infty$.

For $d\ge1$ let $I_d$ be the indicator of $[0,2^{-d})$ on $\T$. Under Lebesgue measure, $I_d(2^b x)=1$ means binary positions $b+1,\ldots,b+d$ of $x$ are all zero. Ignore the null set of ambiguous binary expansions. Define the raw nonnegative block
\[
 Z_s(x)=\sum_{q=1}^{L_s} I_{d_s}(2^{q d_s}x).
 \tag{5}\label{eq:raw}
\]
The digit windows are disjoint, so $Z_s$ is binomial with mean
\[
 \lambda_s=L_s2^{-d_s}=1/t_s\le1.
\]
For a trial as in \eqref{eq:exponents}, let success mean that for at least one $j\in\{\ell,\ldots,L_s\}$, the $d_s$ binary digits beginning just after position $M+jd_s$ are all zero. The $L_s-\ell+1$ candidate windows of a trial are disjoint. If the trial succeeds, for every $r=0,\ldots,\ell-1$ one of the summands in $Z_s(2^{M+r d_s}x)$ uses that same window, by taking $q=j-r\in\{1,\ldots,L_s\}$. Hence the block value is at least one at every dilate in this trial.

A single trial succeeds with probability
\[
 1-(1-2^{-d_s})^{L_s-\ell+1}
 \ge 1-\exp(-\lambda_s/2)\ge \lambda_s/4 .
\]
The choice of $M$ makes \emph{all} trial success events independent, including across stages. A stage is successful if some trial succeeds. Its probability is at least $1-\exp(-t_s\lambda_s/4)=1-e^{-1/4}>0$, and stage-success events are independent. No assertion of independence between the \emph{functions} $Z_s$ is needed.

\subsection*{Choosing the function}
Fix $p>1$ and $g$. Write $\overline g(n)=\max_{1\le m\le n}g(m)$. If $g=o(E_p)$, then $\overline g=o(E_p)$ because $E_p$ is increasing and unbounded: the maximum over any initial finite set is bounded, and for all sufficiently large $m$ one has $g(m)\le\varepsilon E_p(m)$. In particular
\[
 \liminf_{s\to\infty}\frac{\overline g(N_s)}{t_s^{1/p}}=0 .
 \tag{6}\label{eq:envelope}
\]
If $g$ is nondecreasing and only $\liminf g/E_p=0$ is assumed, choose arbitrarily large $u$ witnessing small $g(u)/E_p(u)$ and take $N_s\le u<N_{s+1}$. Then $g(N_s)\le g(u)$ and
\[
 \log(N_s+1)=(2t_s-1)\log 2
 \ge \tfrac13(4t_s-1)\log 2
 \ge \tfrac13\log(u+1).
\]
Since $\overline g=g$ in this case and $\log(N_s+1)$ is comparable to $t_s$, \eqref{eq:envelope} holds again.

Take increasing stage indices $s_k$ and amplitudes
\[
 a_k=2^{-k-2}, \qquad B_k=a_k t_{s_k}^{1/p}
\]
such that $B_k\ge4$ and $B_k/4\ge k\overline g(N_{s_k})$. This is possible for each $k$ by \eqref{eq:envelope} and $t_s\to\infty$. Set
\[
 G=\sum_{k\ge1} B_k Z_{s_k},\qquad f=G-\int_\T G .
\]
For a binomial random variable $Z$ of mean $\lambda\le1$, the elementary bound $x^p\le C_p(e^x-1)$ for nonnegative integers $x$ and the binomial moment-generating function give
$\E Z^p\le C_p(\exp((e-1)\lambda)-1)\le C_p'\lambda$.
Thus $\|B_k Z_{s_k}\|_p\le C_pB_k t_{s_k}^{-1/p}=C_p a_k$, summable. Minkowski and monotone convergence give a well-defined nonnegative $G$ finite almost everywhere and in $L^p$. Centering subtracts a finite constant
\[
 \int_\T G=\sum_k B_k/t_{s_k}
   =\sum_k a_k t_{s_k}^{1/p-1}\le\sum_k a_k<1.
\]
Changing $G$ on its null set does not affect the almost-everywhere statement, since all dilations preserve measure and only countably many are used.

If a selected stage $s_k$ succeeds, consider the endpoint $N=P+\ell=2P+1$ of a successful trial. At least $\ell=P+1>N/2$ of the first $N$ terms have an active contribution $\ge B_k$; all other uncentered blocks and terms are nonnegative. Hence $A_N(f;x)\ge B_k/2-1\ge B_k/4\ge k g(N)$. Stage-success events are independent and have probabilities bounded below, so along the deterministic subsequence $(s_k)$ infinitely many occur almost surely by second Borel--Cantelli. The trial endpoints tend to infinity. This proves \eqref{eq:lower} and the corollary.

\section{Scope and attribution}
The example question in the catalogue is not newly disproved here: Ho's theorem already gives its negative answer. The upper endpoint and the failure of every smaller gauge, on one fixed sequence, sharpen the growth-estimation part and answer Ho's specific order-without-loss question in the universal-gauge sense. There is no assertion of a positive limiting constant for $\sqrt{\log N}$; the upper estimate says that every individual normalized average converges to zero at that scale. The result does not classify sequences with bounded ratios or functions under additional Fourier assumptions.

\begin{thebibliography}{9}
\bibitem{erdos} P. Erd\H{o}s, Problems and results on diophantine approximations, \emph{Compositio Math.} 16 (1964), 52--65. \href{https://www.numdam.org/item/CM_1964__16__52_0/}{Numdam archive}.
\bibitem{ho} B. S. Ho, Counterexamples for lacunary dilates via dyadic spike blocks, \href{https://arxiv.org/abs/2604.18535v3}{arXiv:2604.18535v3}, 23 October 2026.
\end{thebibliography}
\end{document}
